% Source printed pp. 507-508 / PDF pp. 515-516.
\subsection{\texorpdfstring{$K$-半稳定 Fano 簇的模空间}{K-半稳定 Fano 簇的模空间}}\label{src:10.1.4-subsection}
在按 Kodaira 维数的分类中, Fano 簇与一般型簇位于相反的一端: 它们的反典范线丛丰沛. Fano 簇的复几何研究有着丰富的历史, 以 K\"ahler--Einstein 问题为中心: 寻找 K\"ahler 度量 $\omega$, 满足 $\operatorname{Ric}(\omega)=\lambda\omega$, 其中 $\lambda$ 为常数. 在 K\"ahler, Calabi 和 Matsushima 的早期工作基础上, Yau 在二十世纪八十年代猜测, K\"ahler--Einstein 度量的存在应等价于某个代数稳定性条件. Tian 在 \cite{Tia97} 中提出了 $K$-稳定性的概念, Donaldson 则发现这一概念确实是代数性的 \cite{Don02}. 由此产生了著名的 Yau--Tian--Donaldson 猜想: Fano 簇具有 K\"ahler--Einstein 度量当且仅当它 $K$-多重稳定. 这一猜想在 \cite{CDS15a,CDS15b,CDS15c,Szé16} 中得到解决.

在 $2010$ 年以前, 由于存在``等平凡退化'', 例如 $\mathbb P^n\times\mathbb A^1$ 中光滑二次超曲面 $x_0^2+\cdots+x_{n-1}^2+tx_n^2=0$ 退化为锥, 人们几乎不认为能用分离的代数模空间分类 Fano 簇. 然而, Xu 及其合作者在过去十年中取得了惊人的进展, 最终成功构造了 $K$-半稳定 Fano 簇的代数模空间. $K$-半稳定性通过极化簇 $(X,L)$ 的\newterm{检验构形}{test configuration}来定义: 它是 $\mathbb G_m$-等变平坦族 $\mathcal X\to\mathbb A^1$, 带有 $\mathbb G_m$-等变相对甚丰沛线丛 $\mathcal L$, 满足 $1\in\mathbb A^1$ 上的纤维对某个 $r>0$ 同构于 $(X,L^{\otimes r})$. 每个检验构形可视为从 $\Theta=[\mathbb A^1/\mathbb G_m]$ 到极化簇叠的映射 (见命题~\bookref{7.9.7}). 给定检验构形, 可写
\[
h^0(X,L^{\otimes k})=a_0k^n+a_1k^{n-1}+\cdots,
\]
其中 $a_0=L^n/n!$; 若假设 $X$ 正规, 则由 Riemann--Roch 有 $a_1=-(L^{n-1}\cdot K_X)/(2(n-1)!)$. 在 $\mathrm H^0(X,L^{\otimes k})=\mathrm H^0(\mathcal X_0,\mathcal L_0^{\otimes k/r})$ 上有 $\mathbb G_m$-作用, 其权可写成 $w_k=b_0k^{n+1}+b_1k^n+\cdots$. 将 $(\mathcal X,\mathcal L)$ 与平凡族 $(X\times\mathbb A^1,p_1^*L)$ 粘合, 延拓为 $\mathbb P^1$ 上的 $\mathbb G_m$-等变族 $(\widetilde{\mathcal X},\widetilde{\mathcal L})$. 由等变 Riemann--Roch 有 $b_0=(\overline{\mathcal L}^{1/r})^{n+1}/(n+1)!$; 若假设 $\mathcal X$ 正规, 则 $b_1=-(\overline{\mathcal L}^{1/r})^n\cdot K_{\overline{\mathcal X}/\mathbb P^1}/(2n!)$. $(\mathcal X,\mathcal L)$ 的 \newterm{Futaki 不变量}{Futaki invariant}定义为
\[
\operatorname{Fut}(\mathcal X,\mathcal L):=\frac{a_1b_0-a_0b_1}{a_0^2}.
\]
若对每个检验构形 $(\mathcal X,\mathcal L)$ 都有 $\operatorname{Fut}(\mathcal X,\mathcal L)\geq0$, 则极化簇 $(X,L)$ 称为\newterm{$K$-半稳定}{$K$-semistable}; 若它 $K$-半稳定, 且 $\operatorname{Fut}(\mathcal X,\mathcal L)=0$ 仅对乘积检验构形成立, 则称为\newterm{$K$-多重稳定}{$K$-polystable}.

这个定义初次读来可能令人困惑, 但它与用 Hilbert 和 Chow 概形研究 GIT 稳定性时的 Hilbert--Mumford 指标密切相关 (见 \S\bookref{6.8.2-subsection}). 现在限制于 $X$ 是 Fano 簇且 $L=-K_X$ 的情形. 此时 $X$ 必然正规. 为了得到性质良好的模问题, 需限制在温和的奇点类别中, 这里是 \newterm{Kawamata 对数终端奇点}{Kawamata log terminal singularities} (klt), 即正规且差异严格大于 $-1$ (比 slc 的差异 $\geq-1$ 的条件稍强). 记 ${\mathcal F\mathit{anos}}_{n,V}$ 为参数化满足 Koll\'ar 条件的维数 $n$, 体积 $V$ 的 $\mathbb Q$-Fano klt 簇族 $X\to S$ 的代数叠. Fano 簇 $(X,-K_X)$ 为 $K$-半稳定, 当且仅当 $(X,-K_X)\in{\mathcal F\mathit{anos}}_{n,V}$ 相对于 CM 线丛是 $\Theta$-半稳定的 (采用注记~\bookref{8.4.8} 中的术语).

与 KSB 模空间一样, 模空间 $M_{n,V}^{K\text{-}\mathrm{ps}}$ 的六步构造中每一步都极为困难, 也是 $K$-稳定性纲领的重要成果. $K$-半稳定性的开放性由 \cite{BLX22} 建立. 这说明 $K$-半稳定 Fano 簇的轨迹 ${\mathcal F\mathit{anos}}_{n,V}^{K\text{-}\mathrm{ss}}$ 是 ${\mathcal F\mathit{anos}}_{n,V}$ 的开子叠, 因而是代数叠. $K$-半稳定性的有界性由 \cite{Jia20} 证明. 模空间 $M_{n,V}^{K\text{-}\mathrm{ps}}$ 的分离性由 \cite{BX19} 证明. \cite{ABHLX20} 将这一结果加强为 ${\mathcal F\mathit{anos}}_{n,V}^{K\text{-}\mathrm{ss}}$ 是 $\Theta$-完备且 S-完备的; 这既通过存在定理 (\bookref{7.10.1}) 给出分离模空间的存在性, 也蕴含 $K$-多重稳定 Fano 簇的自同构群约化. $M_{n,V}^{K\text{-}\mathrm{ps}}$ 的固有性由 \cite{BHLLX21,LXZ22} 证明, 射影性则由 \cite{CP21b,XZ20b} 证明. 还有到对数 Fano 对的推广. 完整论述见 \cite{Xu25}.
